\section{Formalism}
\label{sec:formalism}

This article is a critical integrative synthesis and operational proposal, not
a systematic review. A project registry records query strata, primary-source
checks, source notes, competing explanations, falsification criteria, and a
claim--evidence matrix; searches were expanded adversarially whenever a field
offered a likely redundancy account. Claims are labelled as established field
results, project evidence, synthesis inferences, or ICP hypotheses, and project
papers are never counted as independent replication.

\subsection{Indexed future-law families}

For a candidate retained state or record $m$, intervention family $\Pi$,
observable algebra $\mathcal A$, boundary and resolution declaration $B$,
resource declaration $C$, and horizon $\tau$, define
\begin{equation}
 \Efamily(m;\Pi,\mathcal A,B,C)
 =\left\{P^{\pi}_{\mathcal A}
   \bigl(F_{t:t+\tau}\mid m,B,C\bigr):\pi\in\Pi\right\}.
 \label{eq:future-family}
\end{equation}
The corresponding reachable support for policy $\pi$ is
\begin{equation}
 \Rset^{\pi}_{\tau}(m)=
 \operatorname{supp}P^{\pi}_{\mathcal A}
 \bigl(F_{t:t+\tau}\mid m,B,C\bigr),
 \label{eq:reachable-support}
\end{equation}
with topology and tolerance specified for continuous spaces.

Equation~\eqref{eq:future-family} is not itself an ICP law; it is shared
comparison language. A \emph{constraint} requires a matched contrast between
families, for example a history-formed carrier $m$ and a reference $m_0$
prepared with the same declared ordinary state variables and resources:
\begin{itemize}
  \item a \textbf{hard constraint} changes reachable support on a common path
  or outcome space;
  \item a \textbf{soft constraint} changes weights while support may remain
  fixed; and
  \item an \textbf{operational constraint} changes a reachability or
  convertibility preorder under the same allowed operations.
\end{itemize}

\subsection{A relative-entropy constraint diagnostic}

For a declared comparison $m$ versus $m_0$ at $x$, one distribution-sensitive
member of the future-family comparison is
\begin{equation}
 \mathfrak C_{x,\tau}^{\pi}
 (m\Vert m_0;\mathcal A,B,C)
 :=D_{\mathrm{KL}}\!\left(
 P^{\pi}_{\mathcal A}(F_{t:t+\tau}\mid m,B,C)
 \middle\Vert
 P^{\pi}_{\mathcal A}(F_{t:t+\tau}\mid m_0,B,C)
 \right),
 \label{eq:icp-constraint-divergence}
\end{equation}
provided both laws use the same outcome space, intervention, horizon, measure,
and coarse-graining, with the first law absolutely continuous with respect to
the second \citep{KullbackLeibler1951Information}. This is an indexed
diagnostic, not the definition of ICP: the full family retains outcome
identity, topology, interventions, and operational order that a single scalar
discards. For a soft constraint it is the natural quantity; for a hard
constraint it has an exact support specialization. If the post-intervention
law is uniform on $\Rset_x^{\mathrm{post}}\subseteq\Rset_x^{\mathrm{prior}}$
and the reference law is uniform on $\Rset_x^{\mathrm{prior}}$, then
\begin{equation}
 D_{\mathrm{KL}}(U_x^{\mathrm{post}}\Vert U_x^{\mathrm{prior}})
 =\log\frac{|\Rset_x^{\mathrm{prior}}|}{|\Rset_x^{\mathrm{post}}|}.
 \label{eq:iota-kl-identity}
\end{equation}
The direction matters: the reversed divergence is infinite under strict
support pruning.

\subsection{The adjacent frontier and signed possibility-space flux}

The adjacent possible becomes operational only after its outcome types and
notion of ``one step'' are declared. Let $r_\ell:\Omega\rightarrow\mathcal
K_\ell$ be a frozen scale-$\ell$ map from micro-outcomes to registered types or
roles, and let $\mathcal V_t^\ell$ be the types actualized by time $t$. Define
\begin{equation}
 \mathcal J_{t,\tau}^{\pi,\ell}(m_t)
 :=r_\ell\!\left(\Rset_\tau^\pi(m_t)\right)\setminus\mathcal V_t^\ell .
 \label{eq:adjacent-frontier}
\end{equation}
When $\tau$ is one registered transition, this is the novelty-bearing sector of
Kauffman's adjacent possible: types not yet actual but reachable in one step
from the actual \citep[p.~142]{Kauffman2000Investigations}. It depends on
$r_\ell,\pi,\tau,B,$ and $C$; changing a classifier, resource budget, or policy
can change the set.

For matched pre/post supports $R_0$ and $R_1$, retain the identity of the
change rather than only its cardinality:
\begin{equation}
 K=R_0\cap R_1,\qquad E=R_1\setminus R_0,\qquad L=R_0\setminus R_1,
 \label{eq:possibility-flux-ledger}
\end{equation}
the retained, enabled, and lost possibilities. A signed log-volume deformation
is $\chi(R_1:R_0):=\log[\mu(R_0)/\mu(R_1)]$. For nonempty finite uniform nested
supports it gives the exact dual pair
\begin{equation}
 \chi(R_1:R_0)=
 \begin{cases}
 D_{\mathrm{KL}}(U_{R_1}\Vert U_{R_0}),&R_1\subseteq R_0
 \quad\text{(pruning)},\\[3pt]
 -D_{\mathrm{KL}}(U_{R_0}\Vert U_{R_1}),&R_0\subseteq R_1
 \quad\text{(enabling)}.
 \end{cases}
 \label{eq:signed-kl-duality}
\end{equation}
Thus $\chi>0$ denotes net contraction and $\chi<0$ net expansion. For
non-nested supports $\chi$ is only a net-volume contrast and can vanish under
complete identity turnover, so the ledger in
Eq.~\eqref{eq:possibility-flux-ledger} must be reported alongside it.

This duality makes one biologically familiar pattern precise. A retained
carrier may restrict immediate microstate options while enabling more
long-horizon role types,
\begin{equation}
 \chi_{\mathrm{micro},1}>0,\qquad
 \chi_{\mathrm{role},h}<0\quad(h\ge h_\star),
 \label{eq:constraint-enablement-crossover}
\end{equation}
under the same randomized carrier intervention. Enabling-constraint and
scaffolding accounts already anticipate this qualitative pattern
\citep{Kauffman2014Prolegomenon,Griesemer2024Scaffolded}; it becomes an ICP
result only if $h_\star$ is predicted from independently measured carrier
lifetimes and confirmed on held-out systems.

\subsection{The persistent-carrier criterion}
\label{sec:criterion}

A proposed informational carrier $M$ passes the ICP criterion only if four
gates are met:
\begin{description}
  \item[Formation.] Controlled histories create the candidate distinction.
  \item[Sufficiency.] The distinction is sufficient for the declared future-law
  family to a frozen error tolerance.
  \item[Carrier intervention.] Ablating, replacing, transplanting, or scrambling
  the proposed carrier changes a preregistered component of the future family,
  while matched shams and resource controls do not.
  \item[Robustness.] The result survives adequate state augmentation and
  reasonable representation, partition, and coarse-graining changes, or its
  representation dependence is stated as a boundary.
\end{description}
Persistence across carrier succession adds one further demand. Let
$c_0,c_1,\ldots$ be physical realizations at successive times. They instantiate
the same role only relative to a declared equivalence relation
$\sim_{\Pi,\mathcal A,B,C,\tau}$ under which replacement preserves the relevant
future-response family. Material disjointness alone does not establish
sameness of role; sameness of an analyst's label does not establish physical
transport.

\subsection{Novelty and failure tests}

The strongest null is complete-state ordinary dynamics: once every causally
relevant variable, law, input, boundary, and control is included, there is no
additional historical variable. ICP earns distinct content only by delivering
at least one of the following beyond the best native model: a minimal
sufficient physically embodied equivalence class; a representation-stable
carrier-succession invariant; a new impossibility or conversion result; a
demonstrably superior compression of intervention--response families; or a
novel preregistered prediction. Disappearance under adequate state
augmentation, failure of matched carrier interventions, or dependence on one
convenient representation would lower or defeat the relevant claim.
